%! AP Statistics — Unit 1, Lesson 1.3 — Group Activity
\documentclass[10pt]{article}
\usepackage{apstats-article}
\usepackage{apstats-boxes}

%=============================================================================
\begin{document}

\pageheader{Unit 1, Lesson 1.3}{Group Activity: Frequency Tables \& Fair Comparisons}
\namepartnerperiod

\vspace{0.1in}

\begin{tcolorbox}[colback=sky, colframe=navy, boxrule=0.8pt, arc=2mm,
    left=3mm, right=3mm, top=2mm, bottom=2mm]
\small
\textbf{Part 1} (group, 8 min): Use the raw data below to build a frequency and
relative frequency table. Then write a description of the distribution.\\[4pt]
\textbf{Part 2} (group, 5 min): A second dataset of different size is provided.
Compare the two distributions using relative frequency.\\[4pt]
\textbf{Part 3} (individual, 3 min): Reflection question.
\end{tcolorbox}

\vspace{0.14in}

% ── PART 1 ───────────────────────────────────────────────────────────────────
{\large\bfseries\color{navy} Part 1 \quad Build the Table}

\vspace{0.1in}

\noindent\small A survey of \textbf{30 middle school students} asks: ``What is your
favorite sport?'' Results are listed below.

\vspace{0.08in}
\begin{center}
\small
\begin{tabular}{*{6}{p{1.4cm}}}
Soccer   & Basketball & Baseball  & Soccer     & Swimming & Tennis \\
Baseball & Soccer     & Swimming  & Basketball & Soccer   & Soccer \\
Tennis   & Basketball & Soccer    & Swimming   & Baseball & Basketball \\
Soccer   & Tennis     & Basketball& Soccer     & Swimming & Baseball \\
Basketball& Soccer    & Tennis    & Basketball & Soccer   & Baseball \\
\end{tabular}
\end{center}

\vspace{0.14in}

\renewcommand{\arraystretch}{2.0}
\begin{tabularx}{\linewidth}{@{} >{\bfseries}p{0.22\linewidth}
                                    C{0.18\linewidth}
                                    C{0.2\linewidth}
                                    C{0.24\linewidth} @{}}
\rowcolor{sky}
\textbf{Sport} & \textbf{Tally} & \textbf{Frequency} & \textbf{Rel.\ Frequency} \\
\midrule
Baseball   & & & \\
Basketball & & & \\
Soccer     & & & \\
Swimming   & & & \\
Tennis     & & & \\
\midrule
\rowcolor{sky}\textbf{Total} & & \blank{2cm} & \blank{2.5cm} \\
\end{tabularx}

\vspace{0.12in}

\begin{enumerate}[leftmargin=1.5em, itemsep=8pt]
  \item Does your total equal 30? \blank{1cm}
        Does your relative frequency column sum to 1? \blank{1cm}

  \item What is the most popular sport and what percentage of students chose it?\\[5pt]
        \writeline

  \item Write a complete sentence describing the distribution of favorite sport
        for these 30 students. Include the most and least common categories.\\[5pt]
        \writeline\\[1pt]\writeline
\end{enumerate}

\vspace{0.14in}

% ── PART 2 ───────────────────────────────────────────────────────────────────
{\large\bfseries\color{navy} Part 2 \quad Compare Two Groups}

\vspace{0.1in}

\noindent\small A separate survey of \textbf{50 high school students} asks the same
question. The frequency table is partially completed. Fill in the relative frequency
column and then answer the comparison questions.

\vspace{0.1in}

\renewcommand{\arraystretch}{2.0}
\begin{tabularx}{\linewidth}{@{} >{\bfseries}p{0.22\linewidth}
                                    C{0.18\linewidth}
                                    C{0.2\linewidth}
                                    C{0.24\linewidth} @{}}
\rowcolor{sky}
\textbf{Sport} & \textbf{Frequency} & \multicolumn{2}{c}{\textbf{Rel.\ Frequency}} \\
\midrule
Baseball   & 8  & \multicolumn{2}{c}{\blank{3cm}} \\
Basketball & 15 & \multicolumn{2}{c}{\blank{3cm}} \\
Soccer     & 14 & \multicolumn{2}{c}{\blank{3cm}} \\
Swimming   & 7  & \multicolumn{2}{c}{\blank{3cm}} \\
Tennis     & 6  & \multicolumn{2}{c}{\blank{3cm}} \\
\midrule
\rowcolor{sky}\textbf{Total} & 50 & \multicolumn{2}{c}{\blank{3cm}} \\
\end{tabularx}

\vspace{0.12in}

\begin{enumerate}[leftmargin=1.5em, itemsep=8pt]

  \item A student says: ``Soccer is the most popular sport in both groups because
        both tables show soccer near the top.'' Is this a valid comparison?
        Why or why not?\\[5pt]
        \writeline\\[1pt]\writeline

  \item \textbf{Using relative frequencies}, which group --- middle school or high
        school --- shows a stronger preference for basketball? Show your evidence.\\[5pt]
        \writeline\\[1pt]\writeline

  \item Suppose a third school has 200 students and 60 choose soccer. Is soccer
        more popular there than in the middle school group? Explain using relative
        frequency.\\[5pt]
        \writeline\\[1pt]\writeline

\end{enumerate}

\vspace{0.14in}

% ── PART 3 ───────────────────────────────────────────────────────────────────
{\large\bfseries\color{navy} Part 3 \quad Individual Reflection}

\vspace{0.1in}

\begin{tcolorbox}[colback=goldbg, colframe=goldacc, boxrule=0.8pt, arc=2mm,
    left=3mm, right=3mm, top=2mm, bottom=2mm]
\small\textbf{In your own words:} Why is a relative frequency table more useful than
a frequency table when comparing two groups? Give a specific example from today's
activity to support your answer.\\[6pt]
\writeline\\[1pt]\writeline\\[1pt]\writeline
\end{tcolorbox}

\end{document}
